\section{Implementation in Mesa}
\label{sec:implementation}

EDEM and DE are implemented in the Python agent-based modelling
framework Mesa~\citep{masad2015mesa}, replacing the earlier NetLogo
prototypes archived at
\url{https://github.com/sibmike/netlogo-realestate-simulations}. The
choice of Mesa is pragmatic: its turtle--patch primitives map cleanly
onto NetLogo's, but the surrounding Python ecosystem (NumPy, pandas,
matplotlib, pytest) supports batched experimentation, deterministic
seeding, and reproducible figure regeneration that the original
NetLogo code cannot. The Mesa port and all experiments described
herein are released under MIT licence at
\url{https://github.com/sibmike/realestate-edem-python}; the
companion paper sources (CC-BY-4.0) live in the same repository.

\subsection{Spatial structure and agents}
\label{sec:implementation-space}

The market is staged on \texttt{mesa.space.MultiGrid(32, 32, torus=True)},
matching the wrapped 32$\times$32 patch world of the original
NetLogo. Per-cell state ($v_{t}(h)$, $v^{*}(h)$, last sale price,
last sale tick) is stored as a lightweight \texttt{Home} dataclass
attached to each grid cell, not as a Mesa Agent --- treating the 1024
cells as Agents would multiply scheduler overhead without benefit,
since cells never schedule behaviour.

Buyers and sellers are Mesa Agents:
\begin{itemize}[leftmargin=*,nosep]
\item \texttt{Seller} is immobile (one per home), holds an ask price
      $a_{t}(s)$ and a patience timer, and accumulates incoming bids
      until the timer elapses.
\item \texttt{Buyer} is mobile; each tick it executes a NetLogo-style
      \emph{wiggle} (uniform heading change in $[-90^{\circ}, +90^{\circ}]$)
      followed by a unit forward step.
\end{itemize}
Bids are NetLogo undirected links in the original; we represent them
as mirrored entries in Python dictionaries on each side (each agent's
\texttt{bids} dict is keyed by the counterparty). The mirror is
maintained in a single helper to avoid drift.

\subsection{Market price and clearing}
\label{sec:implementation-clearing}

For DE, the rolling 25-sale market price of \cref{eq:price-rolling}
is implemented as a \texttt{collections.deque(maxlen=25)} owned by
the model. Sales are committed via a single \texttt{complete\_sale}
method that records the sale, advances the deque, removes both
counterparties from the grid, and triggers the balancer to spawn a
replacement pair (per the NetLogo \texttt{add\_seller} /
\texttt{add\_buyer} idiom).

For the EDEM speculative-market variant, no sales are recorded.
Instead the model maintains a per-tick \texttt{cycle\_counter}
initialised to $\texttt{init\_patience} - 1$. When the counter
reaches zero, the model fires \cref{eq:value-update} on every home in
a single pass and resets the counter. The off-by-one initialisation is
deliberate: it ensures the value update reads buyers' \texttt{is\_yellow}
flags one tick before the buyers' patience timers reset them, matching
the order of operations in the NetLogo \texttt{to go} procedure.

\subsection{Balancer}
\label{sec:implementation-balancer}

DE balances populations every $T_{B}$ ticks by recomputing the linear
target $\hat{Q}_{s}(p), \hat{Q}_{d}(p)$ at the current market price
and adjusting by one agent in each direction. Spawning a replacement
on every sale (the NetLogo behaviour) plus this slow drift toward the
linear target keeps the population near \mbox{$Q_{s}^{*} = Q_{d}^{*}$}
without introducing artificial discrete jumps.

The EDEM balancer implements \cref{eq:Qs-update,eq:Qd-update} in
finite-population form. The fractional balancer coefficient $C_{b}$ is
realised as $\lfloor |C_{b}| \rfloor$ deterministic swaps per epoch
plus one Bernoulli swap with probability $|C_{b}| - \lfloor |C_{b}|
\rfloor$; the sign of $C_{b}$ controls direction. A population floor
prevents the balancer from killing the last agent on either side; the
1-vs-many state that emerges with strong $|C_{b}|$ is kept as a feature
of the model rather than papered over (\cref{sec:exp-balancer-sweep}).

\subsection{Bid-acceptance rule (Cond.~2)}
\label{sec:implementation-cond2}

The original \citet{arbuzov2018de} prose stated the buyer's
bid-acceptance rule with the strict inequality
$\beta < \mathrm{mean}(\mathcal{B}_{b})$, while the accompanying
NetLogo source compared with the opposite sign. Empirically the
prose-as-written rule causes a runaway price collapse, since it
selects every sale to clear at the buyer's lowest bid; only the
NetLogo-as-coded rule reproduces the stable equilibrium reported in
the original Run~1. We adopt \cref{eq:cond2} (the NetLogo-as-coded
rule, restated economically) as the canonical Cond.~2 and document
this discrepancy explicitly in the source. The Mesa implementation
exposes both rules as a parameter (\texttt{accept\_rule="netlogo"}
or \texttt{"prose"}) for sensitivity analysis.

\subsection{Determinism, parallelism, and reproducibility}
\label{sec:implementation-repro}

The Mesa \texttt{Model.rng} is seeded explicitly per run, and all
agent-side randomness (epsilon draws, heading wiggles, balancer
Bernoulli) draws from this single stream. Each experiment script in
\path{python_simulation/experiments/} runs $N \ge 8$ independent seeds
of the same parameter vector and stacks the per-tick model reporters
into a single Parquet dataset; figures plot the median across seeds
plus the 10th--90th percentile band. Pytest unit tests
(\path{python_simulation/tests/}) cover the rolling market price, both
Cond.~2 rules, the linear and fractional balancers, the EDEM
epoch-timing trick, and the $C_{b}$-sign inversion.

\subsection{Per-tick pseudocode}
\label{sec:implementation-pseudocode}

The complete tick is:

\begin{lstlisting}[language=Python]
def step(self):  # Model.step
    self.agents.shuffle_do("step")     # buyers + sellers, randomised
    self.balancer.step(self)           # DE: every T_B; EDEM: per-epoch
    self.datacollector.collect(self)
\end{lstlisting}

\noindent
where \texttt{Buyer.step} runs \emph{buy + wiggle + move} in NetLogo
order and \texttt{Seller.step} runs the patience-timer logic
(timer decrement; on timeout, either drop the highest bid below ask
and lower the ask, or offer to sign the highest bidder, applying
\cref{eq:cond2}). The full source is under 600 lines including
docstrings.
